% \documentclass[a4paper,11pt,twoside]{article}
\documentclass[a4paper,11pt,oneside]{article}
\usepackage{amsmath,amsthm,amssymb}
\usepackage{enumitem}
\usepackage[latin1]{inputenc}
\usepackage[swedish,english]{babel}
\usepackage[pdftex]{graphicx}
\usepackage{xcolor}
\usepackage[pdftitle={The PDF title}, % *** edit
pdfauthor={The Authors}, % *** edit
pdffitwindow=true,
breaklinks=true,
colorlinks=true,
urlcolor=blue,
linkcolor=red,
citecolor=red,
anchorcolor=red]{hyperref}
\usepackage{doi}
\usepackage[numbers,sort]{natbib}
\usepackage{itkompendium}

\usepackage{listings}

\definecolor{dkgreen}{rgb}{0,0.6,0}
\definecolor{gray}{rgb}{0.5,0.5,0.5}
\definecolor{mauve}{rgb}{0.58,0,0.82}

\lstset{frame=tb,
  language=C,
  aboveskip=3mm,
  belowskip=3mm,
  showstringspaces=false,
  columns=flexible,
  basicstyle={\small\ttfamily},
  numbers=none,
  numberstyle=\tiny\color{gray},
  keywordstyle=\color{blue},
  commentstyle=\color{dkgreen},
  stringstyle=\color{mauve},
  breaklines=true,
  breakatwhitespace=true
  tabsize=3
}

% **************************************************************************

% project report checklist
\newlist{todolist}{itemize}{2}
\setlist[todolist]{label=$\square$}
\newcommand{\itemdone}{\item[$\boxtimes$] }
     
\numberwithin{equation}{section}
\numberwithin{table}{section}
\numberwithin{figure}{section}

% how to reference figures, tables and sections:
\newcommand{\figref}[1]{Figure~\ref{fig:#1}}
\newcommand{\tabref}[1]{Table~\ref{tab:#1}}
\newcommand{\secref}[1]{\S\ref{sec:#1}}

% do/don't
\definecolor{Green}{rgb}{0.15,0.75,0.35}
\definecolor{Red}{rgb}{0.75,0.05,0.1}
\newcommand{\DO}[1]{\textcolor{Green}{\{Do:\}}#1}
\newcommand{\DONT}[1]{\textcolor{Red}{\{Don't:\}}#1}
\newcommand{\blabla}[1]{\textit{\{#1...\}}}

% how to show that you have reviewed something
\newcommand{\review}[1]{\textcolor{blue}{#1}}

% ******************************************************************

\begin{document}

\selectlanguage{english}

\titel{Project in Computational Science Report Template} % *** edit
%\undertitel{Subtitle (if needed)} % (only if needed!) *** edit

\marginaltext{PROJECT REPORT}

\forfattare{Stefan Engblom, Author Two, Author Three} % *** edit
% generally use the stile "Surname Familyname" (+initials if you so
% prefer)

\rapportnummer{Project in Computational Science: Report} 

\datum{\today}

\itkompendium

\begin{abstract} % *** edit
  This is the abstract of the report. An abstract is a short focused
  summary of what you have done. It is important that it can be
  understood in itself. A good abstract attracts a reader!

  \DO{Explain the general problem and context. Focus a bit on what you
    did and what the general outcome and conclusion was. Use a formal
    language.}

  \DONT{Avoid informal language: no use of \emph{don't}, \emph{isn't}
    etc. Never pick variables from your setting and insert into the
    abstract. Avoid references unless you have very good reasons. Do
    not use acronyms without defining them.}
\end{abstract}

\clearpage

\tableofcontents

\clearpage

% ******************************************************************

\section{Introduction}
\label{sec:intro}
\label{sec:intro}

The introduction should give the reader the neccessary background. For
example, start with a bird's perspective and zoom in on the smaller
details investigated in the present project.

Citations supporting your statements look like this:
\cite{artref,procref}, that is, a cited article and a conference
proceeding. Single book citations \cite{bookref} or master's theses
\cite{thesisref} may also be used. Frequently, so-called
\emph{preprints} are cited before they even have been published
\cite{preprintref}. It goes without saying that one might also want to
cite web-pages, like a specific URL \cite{urlref} or, perhaps more
common, a Wikipedia-page \cite{wikiref}.

\DO{Remember that the citations are part of your sentence and should
  be included in the sentence as such, i.e., come before a `.' at the
  end of the sentence, are separated by commas where appropriate,
  etc.}  \DO{Look at the Bibliography file and/or check any Bib\TeX\
  warnings and make sure to fill out all the appropriate fields for
  the type of reference you use.}

\DONT{Citations should be used to support and reference what you have
  written. Readers should understand the main point you have taken
  from the content of your citation without having to read it
  themselves.}

\DO{ It is common to end the introduction with an overview of the
  structure of the report.} In this example, theory is summarized in
\secref{theory}, and the required methods in \secref{methods}. The
results are to be found in \secref{results} and a discussion around
the central themes of the project is finally offered in
\secref{conclusions}. Note the use of \verb!\secref{<section label>}!
to reference sections.

% ******************************************************************

\section{Theory}
\label{sec:theory}

\blabla{The main theory stems from the observation that} \DO{All
  sections should start with a short introductory sentence or
  paragraph explaining what the section will provide.} Please consider
to select a better section heading than ``theory'' which is rather
generic!

\subsection{Details}
\DO{It is common to use subsections for the separate pieces that
  together make up the section.}

When you describe the details of your work, you might sometimes need
to include an equation. If you, for example, want to describe the
Pythagorean Theorem you would state
\begin{align}
  \label{eqn:pythagoras}
  a^2 + b^2 &= c^2,
\end{align}
with $a$ and $b$ the catheti, and $c$ the hypotenuse of a right-angled
triangle. \DO{Remember to write the equation in such a way that it is
  incorporated in your sentence. A good rule of thumb is that your
  sentence should make sense when read out loud. Also remember to
  describe the unknowns of the equation.}

\DO{When you want to refer to your equation, you need to reference
  properly which in this case is simply ``see
  \eqref{eqn:pythagoras}'', or \emph{rarely}, ``see
  equation~\eqref{eqn:pythagoras}''. Prefer the first shorter
  version.}

\DO{Remember who your reader is going to be. Typically this is a
  another student in the same course and with similar background to
  you before you started the project. So write in such a way that the
  reader can understand what you have done. This means including
  enough detail so others can understand both the background you have
  read about and what you have done.}

\subsubsection{Too small details}
\DONT{It is less common to use `subsubsections' since it is easy to
  loose the flow.}

% ******************************************************************

\section{Methods}
\label{sec:methods}

\blabla{In this section we will introduce the two main methods
  investigated} \DO{Remember to introduce the contents of all
  sections, if however short.}  Please consider to select a better
section heading than ``methods'' which is rather generic!

\subsection{Algorithm solving the main problem}

Consider to take a look at the various \LaTeX\ packages for displaying
algorithms. There are plenty of options available. In \figref{CGfig} a
convergence history is illustrated. Note the use of
\verb!\figref{<figure label>}! to refefence a figure.

\begin{figure}[htp]
  \centering
  \includegraphics[width = 0.35\textwidth,angle = 90]{CGfig}
  \label{fig:CGfig}
  \caption{Conjugated gradient convergence
    illustrated. \blabla{explanation goes here} Figure taken from
    \cite{wikiCGref} (Public domain). \DO{Unless you have made the
      figure yourself, cite all figures and state the permission under
      which you use them.} \DONT{Never use figures without
      permission.}}
\end{figure}

% ******************************************************************

\section{Results}
\label{sec:results}

\blabla{The main results are structured as follows} \DONT{The
  subsections below show how to format figures and tables in your
  text. Don't use these specific separate subsections in your own
  report!}  \DO{Instead, present your results in a way that makes the
  most sense for your project.}

\subsection{Figures}

To generate figures, some tricks may come in handy. For example,
\figref{plot} was generated in Matlab using
\begin{verbatim}
>> xlabel('x [mm]','interpreter','latex')
>> ylabel('y [mm]','interpreter','latex')
>> legend('Method 1','Method 2','location','SE','interpreter','latex')
\end{verbatim}
In order to get nice text in the figure.

Further, a figure can usually be cleaned up a bit by fixing the axis
and additionally scaling it to the final size:
\begin{verbatim}
>> set(gca,'xtick',[0 1],'ytick',1:3)
>> set(gcf,'Position',[100 100 240 180])
\end{verbatim}
This can be done similarly in other programming languages such as
Python and avoids scaling inside \LaTeX\ such that, e.g., the text
looks weirdly scaled.

\begin{figure}[htp]
  \centering
  \includegraphics{plot}
  \label{fig:plot}
  \caption{Comparision of Methods 1 and 2. \blabla{explanation goes
      here} \DO{Remind the reader of the meaning of the important
      variables. Consider to use different markers and not only colors
      for clarity (here: squares and circles, respectively. You should
      also refer to the figure in the main text.}  \DONT{Never insert
      a title in the figure, use the caption instead.}}
\end{figure}

\subsection{Tables and Code}

Tables such as \tabref{results} can also be used. Be sure to also
refer to the table in the text and offer a discussion about the
results. Note the use of \verb!\tabref{<table label>}! to refefence a
table.

Sometimes it can be nice to include code as well, see
Listing~\ref{CodeTest} for an example. However, think carefully of why
the code should be included. Who will actually read and understand it?

\begin{table}[ht]
  \centering % for centered tables
  \begin{tabular}{c r r r}
    Case & Method \#1 & Method \#2 & Method \#3 \\
    \hline
    1 & 50 & 837 & 970  \\
    2 & 47 & 877 & 230  \\
    3 & 31 & -  & -  \\
    \hline
  \end{tabular}
  \caption{Model results. Note that method \#2 and \#3 did not finish
    for the third case. \blabla{more explanation goes here}
    \DO{Explain all notation in the table, e.g., the dashes in this
      example. You should also refer to the table in the main text.}
    \DONT{Professional tables generally \emph{do not have vertical
        lines}.}}
  \label{tab:results}
\end{table} 

\begin{lstlisting}[language = Matlab, label = CodeTest, caption =
  Sample code.]
function c = repeat(a,b)
%REPEAT Repeat elements of vector.
%   C = REPEAT(A,B) returns A repeated by B, see the example
%   below. Both arguments are assumed to be row vectors of the same
%   length and B must contain nonnegative integers only.
%
%   No error-checking is performed.
%
%   Example:
%     a = [0 1 pi inf nan -nan];
%     b = [1 4 0   2   3    0 ];
%     c = repeat(a,b) % returns c = [0 1 1 1 1 inf inf nan nan nan]

% S. Engblom 2005-10-14

% can't beat this!
ix = cumsum(sparse(1,cumsum([1 b]),1));
c = a(ix(1:end-1));
\end{lstlisting}

\subsection{Project report checklist}

\DO{This list should be filled out fully and submitted together
  \emph{with your submissions until the final one!}} Insert it as a
\verb|\subsection*{Project report checklist}| right by the end of
\secref{intro}. In submitting the final and revised version you simply
skip this checklist.

\begin{itemize}
\item Titlepage, Abstract, Table of contents
  \begin{todolist}
  \item The title is the correct and final title
  \item All author's names are there, usually in alphabetic order as
    to the family name
  \item The abstract is informative and focuses on the problem
    approached and what you did
  \item The abstract does not contain undefined abbreviations,
    acronyms, references, equations, etc., unless you have good
    reasons
  \item The table of contents is up to date and makes a structured and
    clear impression
   \itemdone Checked items are marked out using \verb!\itemdone! and
   all checked items are actually taken care of!
  \end{todolist}
  
\item Structure
  \begin{todolist}
  \item The report makes sensible use of sections and subsections:
    avoid subsubsections and appendices unless you have good reasons
  \item In the case there is an appendix, it is explained why it is
    there in its initial paragraph
  \item The first section (introduction) explains what will come in the
    report and in what order (usually in the final paragraph)
  \item Each section starts with a few sentences explaining what will
    come in the section
  \end{todolist}

\item Contents
  \begin{todolist}
  \item Each section is made up of paragraphs which gives a coherent
    flow to a reader and concerns the topic of the section
  \item Each paragraph is made up of sentences which gives a coherent
    flow to a reader and concerns the topic of the paragraph
  \item Other than common knowledge or trivialities, no claims are made
    without stating the reasons or inserting a proper citation which
    supports the claim
  \item Equations are used when needed and are referenced when useful
  \item The report has been consistently spell-checked and proofread by
    all members in the project group
  \item You have received feedback from your supervisor(s) and you have
    made a serious attempt to adress all comments received
  \item Volume: the normal range is about 12--24 pages; in case the
    report is shorter, then there are reasons for this, and in case the
    report is substantially longer, then there are reasons for this
  \end{todolist}
  
\item References, Figures, Tables
  \begin{todolist}
  \item References are used throughout to support claims and direct
    the reader to useful sources of information
  \item The actual items in the reference list by the end uses
    \verb!\doi{}!, in some cases \verb!\url{}!, but \emph{never
      both}. \emph{Never} use references which involves a
    Google-search or similar. Similarly \emph{never} use references
    from ScienceDirect or similar pages as these URLs are not original
    ones!
  \item The use of ISSN/ISBN numbers have been avoided
    altogether. Also, generally, only years are listed, so not
    individual months.
  \item For all cited arXiv preprints, these have been checked in case
    they are actually published. In case they were published, this
    version is rahter the one cited.
  \item No figure is taken from other sources without either a written
    permission (stated) or under an open license (also stated)
  \item Captions for figures and tables explain what information is
    communicated
  \item Captions are a full short paragraph and hence ends with a
    period
  \item Figure titles are avoided; units are clearly desiplayed and
    the font for the axis is reasonably large
  \item All figures included in the report are meaningful and are
    meant to be there since they communicate information
  \item All figures and tables are discussed in the main text and are
    then also referenced
  \end{todolist}

\item Use of AI
  \begin{todolist}
  \item All AI-generated figures are marked out/cited as such, are
    included because they are meaningful, and accurately convey the
    information they are meant to
  \item The report generally and transparently details all uses of
    AI-tools, e.g., as a minium by explaining how it was used and to
    what extend (see the acknowledgment for an example)
  \item Text or code proposed or modified by an AI has been checked
    for correctness and that the meaning is as intended
  \item All members of the project team take full responsibility for
    the whole text
  \end{todolist}
\end{itemize}

\subsection{How to incorporate feedback}

During the course you will get feedback several times. \DO{Be sure to
  update your report according to \emph{all} comments. Ask if you are
  unsure!}

By the end of the course you will get a fuller feedback of the report
as a whole. After getting this feedback you prepare a hopefully final
version. In order to help in the review process the second time, use a
\verb|\review|-command, which helps the reviewer to quickly locate
updated text. \review{For example, this is what the command does to a
  single sentence}. When submitting the accepted version you can
simply redefine the command as in
\begin{verbatim}
  \newcommand{\review}[1]{\textcolor{black}{#1}}
\end{verbatim}
This makes the text look normal again.


% ******************************************************************

% *** it looks nicer if the final conclusion is not mixed with figures
% belonging to other sections, so sometimes either one of \clearpage
% or \cleardoublepage is needed

\section{Conclusions}
\label{sec:conclusions}

In this section you summarize the results and try to offer ideas on
directions for future investigatons. Try to put the investigations
made in the report in a larger perspective: what was the main outcome?
\DO{Take into consideration how long your report in total is and how
  long the conclusion should be relative to that. If you write
  multiple pages including a lot of detail, 1--2 sentences for the
  conclusion are not enough.}

\subsection*{Acknowledgment}

The acknowledgment section in a project report is where you express
gratitude to those who have supported you, including supervisors,
mentors, colleagues, and anyone who provided assistance or
resources. It's a space to recognize their contributions and express
appreciation for their help in completing the project.

\medskip

\noindent
\texttt{ChatGPT 3.5} was used as a text-editing tool in minor parts of
this text. The Acknowledgment section, however, was generated in its
entirety in this way.

% ******************************************************************

\bibliographystyle{abbrvnat}
\bibliography{references.bib}

% ******************************************************************

\clearpage

\appendix

\section{Not so often and for very good reasons}

\blabla{explanation of what this Appendix brings}

\DO{If you absolutely want an Appendix you should explain what the
  purpose is and also refer to it in the main report. Otherwise you
  can most likely skip it?} \DONT{Never add ``all the figures'' or
  ``additional results and code''. Who should care?}

\section{Most likely not}

\blabla{also need to explain this Appendix}

\end{document}
% ******************************************************************
